% SPDX-License-Identifier: Apache-2.0
\section{Defects in the Sources}
\label{sec:defects}

These were found by reading, because no parser exists for either language.
Each has to be settled while merging, since a merged specification cannot
state two things at once.

\subsection{TxHDL}

Eleven, in \code{spec/language.md}. The line numbers are from the file as
committed.

\begin{itemize}
  \item Line 1: the file begins with the literal text \code{:qODq} before
        its title. A stray editor keystroke was committed, and it rendered
        into the title of the first PDF built from the file.
  \item Line 660 states that a module holds only one sequence.
        \code{module Counter} at line 416 declares two, and
        \code{module DualPortMemory} at line 2012 declares two.
  \item Lines 133 and 141: inside one enumeration, two variants share the
        encoding \code{0x02}. Lines 135 and 140 share \code{0x04}.
  \item Line 1112: an array declared with 256 elements is initialised with
        four and a comment.
  \item Line 740: a timeout requires its else block to end in a value, and
        blocks as expressions are never defined.
  \item Line 2043: a port uses an anonymous bus type inline, and the bus
        section never defines that form.
  \item Line 1196: a clock domain is given a frequency literal, and the
        literal section defines no unit suffix.
  \item There is no grammar. A quick reference lists keywords and
        operators, and never their syntax.
  \item Transactions and structs are both product types, with no stated
        rule for choosing one.
  \item Subtype compatibility says a receiver ignores extra fields, and the
        width of the physical channel for that case is not stated.
\end{itemize}

\subsection{LHDL}

The grammar and the examples disagree in sixteen places. Three are enough on
their own: no interface, no modport and no generic parameter list written in
any of the four LHDL files parses against the grammar in the condensed
specification.

\begin{itemize}
  \item An interface item must be a directed port declaration, and every
        interface in every file writes undirected fields.
  \item A modport takes port declarations, which require a type, and every
        modport writes a direction and a name with no type.
  \item A generic declaration puts the name first, and the examples write
        the keyword first.
  \item A tag definition requires a type, and five examples omit it. One
        writes the tag name with a leading marker the grammar does not
        allow. Another writes a comma the grammar does not allow.
  \item A variable declaration has no initialiser, and the examples give
        one.
  \item The parallel-composition statement has two incompatible
        productions, one per file. Every example uses the first.
  \item The operator table offers two assignment operators, and the grammar
        offers one.
  \item A state transition takes a keyword in the grammar, and one example
        omits it.
  \item The prose names functions one way and the grammar another.
  \item Struct fields are comma separated in the grammar and semicolon
        separated in one example.
  \item A module has no port list and no tag in the grammar, and one
        example gives it both.
  \item Neither the domain block nor a bare binding appears among the
        top-level definitions, and both appear at top level in an example.
  \item An instance declaration requires a parenthesised binding list, and
        an example has none.
  \item One example calls an interface instance as a function, and no rule
        covers that.
\end{itemize}

The LL(1) claim is asserted and never checked. One conflict is visible
without tooling: two top-level productions both begin with an optional
\code{pub}, so a predictive parser cannot choose between them on that token
until the prefix is factored out.

\subsection{Lessons for the unified specification}

Both original source drafts include grammars or syntactic claims that diverge
from their own code listings. This divergence persisted because neither draft
was verified by automated tooling. In the unified specification, every
syntactic construct possesses a corresponding reference example in a standalone
file verified by automated test harnesses. This practice enforces the project
requirement that every specification construct must be accompanied by a
passing automated test.
